% Lösungen Einheit 1 von 3 – Unabhängigkeit prüfen: die Produktregel P(A ∩ B) = P(A) · P(B) an Anteilen oder absoluten Häufigkeiten prüfen (Tafel oder Anzahlen erst beschaffen (zusammenbau v0.1)
\einheitenkopf{Einheit 1 von 3 · Unabhängigkeit prüfen: die Produktregel P(A ∩ B) = P(A) · P(B) an Anteilen oder absoluten Häufigkeiten prüfen (Tafel oder Anzahlen erst beschaffen}

\erg{9}{a) prüfen – die Produktregel anwenden: alle drei Zahlen beschaffen und vergleichen}
\erg{10}{a) alle drei stehen in der Tafel: zwei Ränder und ein Feld – $P(A)$ und $P(B)$ am Rand, $P(A \cap B)$ im Feld \quad b) $P(M \cap F) = \frac{20}{100} = 0{,}2$, $P(M) \cdot P(F) = 0{,}4 \cdot 0{,}5 = 0{,}2$ – gleich, also unabhängig \quad c) 12 Erwachsene spielen Volleyball; $P(E \cap V) = \frac{12}{60} = 0{,}2$, $P(E) \cdot P(V) = 0{,}6 \cdot \frac{1}{3} = 0{,}2$ – gleich, also unabhängig \quad d) Ältere mit Bestehen: $9\,920 - 7\,750 = 2\,170$; $P_{Ä}(S) = \frac{2\,170}{3\,100} = 0{,}7$, $P(S) = \frac{9\,920}{12\,400} = 0{,}8$ – ungleich, also abhängig: Ältere bestehen seltener als der Durchschnitt. \quad e) schwer: $\frac{90}{150} = 0{,}6$; nicht schwer: $\frac{150}{450} \approx 0{,}333$ – verschieden, also abhängig \quad f) $P(A) = \frac{4}{16} = 0{,}25$ (Paare 1-4, 2-3, 3-2, 4-1), $P(B) = \frac{4}{16} = 0{,}25$, $P(A \cap B) = \frac{1}{16}$ (nur 4-1); $0{,}25 \cdot 0{,}25 = 0{,}0625 = \frac{1}{16}$ – unabhängig \quad g) Links steht $P_K(O) = 0{,}4$, der Anteil der Onlinebesteller unter den Karteninhabern; rechts $P(O) = 0{,}7 \cdot 0{,}4 + 0{,}15 = 0{,}43$, der Anteil der Onlinebesteller unter allen Kunden. Bedingter und unbedingter Anteil sind verschieden, also sind K und O abhängig.}

10c)

\vierfeldertafel{E}{V}{12,8,20,24,16,40,36,24,60}

\erg{11}{a) Fehler: Lars verwechselt Unabhängigkeit mit Unvereinbarkeit; ein nichtleerer Schnitt sagt nichts über Abhängigkeit. Richtig: $P(M) \cdot P(F) = 0{,}4 \cdot 0{,}5 = 0{,}2 = P(M \cap F)$, also unabhängig. \quad b) Der bedingte Anteil ist der Quotient aus Schnitt und $P(A)$; teilt man die Produktregel durch $P(A)$, steht genau $P_A(B) = P(B)$ da – beide Gleichungen sagen dasselbe, solange $P(A)$ nicht null ist. \quad c) $P(\text{Öko} \cap \text{E}) = \frac{3\,600}{40\,000} = 0{,}09$, $P(\text{Öko}) \cdot P(\text{E}) = 0{,}3 \cdot 0{,}2 = 0{,}06$ – abhängig: Unter Elektroautobesitzern ist der Ökostromanteil $0{,}45$ statt $0{,}3$; sie sind schon überdurchschnittlich Kunden, die Werbung erreicht dort weniger Neukunden. \quad d) $0{,}4$; $0{,}6$; $0{,}5$; $0{,}5$ (erste Stufe M und nicht M; zweite Stufe F und nicht F, in beiden Zweigen gleich) – gleiche Äste, also unabhängig}

11d)

\baumzwei{M/0{,}4,\overline{M}/0{,}6}{F/0{,}5,\overline{F}/0{,}5}{F/0{,}5,\overline{F}/0{,}5}

